{\actuality}
\fixme{Вычислительные системы -- это термин специальности 2.3.2}
\fixme{Удалить ссылки, картинки, связать с sec:ch1}

Процесорные технологии развиваются, предлагая всё более мощные и эффективные вычислительные системы.
По закону Мура~\cite{moore} рост числа транзисторов на кристалле удваивается каждые два года.
Но в последнее время рост числа транзисторов не коррелировал с ростом тактовой частоты, как было ранее~\cite{trend}.
На рисунке~\ref{fig:trend} показано изменение количества транзисторов, количества вычислительных ядер, тактовой частоты вычислительного ядра и энергопотребления.

\begin{figure}[ht]
    \centerfloat{
      \includegraphics[scale=2]{trend.jpg}
    }
    \caption{Изменение количества транзисторов, вычислительных ядер, тактовой частоты и энергопотребления.}
    \label{fig:trend}
\end{figure}

Рост тактовой частоты, который ранее являлся основным драйвером повышения производительности, замедлился.
По пути повышения быстродействия вычислительное ядро приблизилось к своим физическим и технологическим пределам~\cite{physLimits}.
Для дальнейшего роста производительности, инженеры обратились к другим стратегиям, и ключевым направлением развития стала многоядерность.
Это позволило выполнять несколько задач одновременно, повышая общую производительность.
Вследствие этого, сложность систем значительно увеличилась.
В них применяются многоуровневые иерархические архитектуры, параллелизм исполнения требует сложных механизмов синхронизаций, усложняя взаимодействие между различными компонентами~\cite{hieracy}.
Рост сложности современных вычислительных систем опережает развитие языковых средств их описания.
Несмотря на широкое распространение в промышленности языков описания аппаратуры (ЯОА), таких как VHDL, Verilog и его расширенная версия SystemVerilog, процессы верификации, анализа и отладки создаваемых на них моделей остаются крайне трудоемкими.
Отчасти данную проблему решает высокоуровневые языки описания архитектуры (High-Level Synthesis, HLS), которые предоставляют разработчикам современные парадигмы проектирования, более высокий уровень абстракции, а также встроенные возможности для анализа и валидации кода.
Однако, даже с их использованием, обеспечение полной верификации создаваемой аппаратуры на соответствие требованиям технического задания представляется сложной задачей ввиду общей сложности современных вычислительных систем.

Верификация синтеза современных вычислительных систем  критически важна, поскольку она обеспечивает гарантию корректной работы, повышение надёжности и безотказности.
Кроме того, своевременное обнаружение проблем и ошибок позволяет существенно сократить временные затраты на разработку.
В связи с этим, актуальной задачей представляется развитие \fixme{точно развитие???} методов верификации и тестирования программ, описывающих вычислительные системы.

{\aim} данной работы является создание методов и разработка алгоритмов анализа консистентности памяти при исполнении многопоточных программ в многопроцессорных системах с учётом моделей памяти.

{\progress}
\fixme{where???}Здесь про прогресс, litmus, manovit, intell, etc.


Для~достижения поставленной цели необходимо было решить следующие {\tasks}:
\begin{enumerate}[beginpenalty=10000] % https://tex.stackexchange.com/a/476052/104425
  \item построение частично упорядоченного множества операций с памятью в многопоточной программе с учётом модели памяти;
  \item разработка алгоритмов построения всех линейных порядков частично упорядоченного множества, учитывающих свойства эквивалентных перестановок;
  \item разработка компактного и полного описания результатов исполнения многопоточной программы;
  \item разработка алгоритмов анализа результата исполнения многопоточной программы для верификации соответствия аппаратуры заявленной модели памяти;
  \item апробация \fixme{дополнить}.
\end{enumerate}


% Объект исследования представляет собой знание, порождающее проблемную ситуацию
{\object} являются модели памяти многопроцессорных вычислительных систем.\fixme{улучшить}

% Предметом исследования являются технические решения, методы, алгоритмы, технологии для
{\subject} являются модели представления многопоточных программ с учётом модели памяти и алгоритмы их анализа и верификации результата исполнения.

{\novelty}
\begin{enumerate}[beginpenalty=10000] % https://tex.stackexchange.com/a/476052/104425
  \item Разработаны алгоритмы построения линейных порядков частично упорядоченного множества с учётом зависимости и с учётом независимости.
  \item Разработано описание исполнения программы (граф состояний), учитывающее не только численный результат исполнения, но и порядок исполнения, приведшего к численному результату.
  \item Продемонстрирована возможноcть исчерпывающей верификации многопоточной программы с пересекающимися доступами в память в наиболее общем случае.
\end{enumerate}

{\influence}. Теоретическая значимость полученных научных результатов состоит в развитии алгоритмического описания многопоточных программ с учётом моделей памяти.
Практическая значимость научных результатов состоит в разработке инструментария, направленного на выявление критических ошибок многопоточных систем на ранних этапах разработки, что способствует повышению отказоустойчивости и стабильности вычислительных систем.

{\methods} Для решения поставленных задач в работе используются методы комбинаторного анализа, теории графов. \fixme{дополнить}

{\defpositions}
\begin{enumerate}[beginpenalty=10000] % https://tex.stackexchange.com/a/476052/104425
  \item Алгоритмы построения линейных порядков частично упорядоченного множества с учётом зависимости и с учётом независимости.
  \item Модель описания \fixme{множественное число, когда про StateGraph} исполнения многопоточной программы (граф состояний).
  \item Алгоритм верификации результата исполнения многопоточной программы с учётом моделей памяти.
\end{enumerate}

{\reliability} полученных результатов обеспечивается математическими свойствами представления исполнения многопоточной программы. \ Результаты находятся в соответствии с результатами, полученными другими авторами.\fixme{подумать}


{\probation}
Основные результаты работы докладывались~на:
перечисление основных конференций, симпозиумов и~т.\:п.

{\contribution} Автор принимал активное участие \ldots

\ifnumequal{\value{bibliosel}}{0}
{%%% Встроенная реализация с загрузкой файла через движок bibtex8. (При желании, внутри можно использовать обычные ссылки, наподобие `\cite{vakbib1,vakbib2}`).
    {\publications} Основные результаты по теме диссертации изложены
    в~XX~печатных изданиях,
    X из которых изданы в журналах, рекомендованных ВАК,
    X "--- в тезисах докладов.
}%
{%%% Реализация пакетом biblatex через движок biber
    \begin{refsection}[bl-author, bl-registered]
        % Это refsection=1.
        % Процитированные здесь работы:
        %  * подсчитываются, для автоматического составления фразы "Основные результаты ..."
        %  * попадают в авторскую библиографию, при usefootcite==0 и стиле `\insertbiblioauthor` или `\insertbiblioauthorgrouped`
        %  * нумеруются там в зависимости от порядка команд `\printbibliography` в этом разделе.
        %  * при использовании `\insertbiblioauthorgrouped`, порядок команд `\printbibliography` в нём должен быть тем же (см. biblio/biblatex.tex)
        %
        % Невидимый библиографический список для подсчёта количества публикаций:
        \phantom{\printbibliography[heading=nobibheading, section=1, env=countauthorvak,          keyword=biblioauthorvak]%
        \printbibliography[heading=nobibheading, section=1, env=countauthorwos,          keyword=biblioauthorwos]%
        \printbibliography[heading=nobibheading, section=1, env=countauthorscopus,       keyword=biblioauthorscopus]%
        \printbibliography[heading=nobibheading, section=1, env=countauthorconf,         keyword=biblioauthorconf]%
        \printbibliography[heading=nobibheading, section=1, env=countauthorother,        keyword=biblioauthorother]%
        \printbibliography[heading=nobibheading, section=1, env=countregistered,         keyword=biblioregistered]%
        \printbibliography[heading=nobibheading, section=1, env=countauthorpatent,       keyword=biblioauthorpatent]%
        \printbibliography[heading=nobibheading, section=1, env=countauthorprogram,      keyword=biblioauthorprogram]%
        \printbibliography[heading=nobibheading, section=1, env=countauthor,             keyword=biblioauthor]%
        \printbibliography[heading=nobibheading, section=1, env=countauthorvakscopuswos, filter=vakscopuswos]%
        \printbibliography[heading=nobibheading, section=1, env=countauthorscopuswos,    filter=scopuswos]}%
        %
        \nocite{*}%
        %
        {\publications} Основные результаты по теме диссертации изложены в~\arabic{citeauthor}~печатных изданиях,
        \arabic{citeauthorvak} из которых изданы в журналах, рекомендованных ВАК%
        \ifnum \value{citeauthorscopuswos}>0%
            , \arabic{citeauthorscopuswos} "--- в~периодических научных журналах, индексируемых Web of~Science и Scopus%
        \fi%
        \ifnum \value{citeauthorconf}>0%
            , \arabic{citeauthorconf} "--- в~тезисах докладов.
        \else%
            .
        \fi%
        \ifnum \value{citeregistered}=1%
            \ifnum \value{citeauthorpatent}=1%
                Зарегистрирован \arabic{citeauthorpatent} патент.
            \fi%
            \ifnum \value{citeauthorprogram}=1%
                Зарегистрирована \arabic{citeauthorprogram} программа для ЭВМ.
            \fi%
        \fi%
        \ifnum \value{citeregistered}>1%
            Зарегистрированы\ %
            \ifnum \value{citeauthorpatent}>0%
            \formbytotal{citeauthorpatent}{патент}{}{а}{}%
            \ifnum \value{citeauthorprogram}=0 . \else \ и~\fi%
            \fi%
            \ifnum \value{citeauthorprogram}>0%
            \formbytotal{citeauthorprogram}{программ}{а}{ы}{} для ЭВМ.
            \fi%
        \fi%
        % К публикациям, в которых излагаются основные научные результаты диссертации на соискание учёной
        % степени, в рецензируемых изданиях приравниваются патенты на изобретения, патенты (свидетельства) на
        % полезную модель, патенты на промышленный образец, патенты на селекционные достижения, свидетельства
        % на программу для электронных вычислительных машин, базу данных, топологию интегральных микросхем,
        % зарегистрированные в установленном порядке.(в ред. Постановления Правительства РФ от 21.04.2016 N 335)
    \end{refsection}%
    \begin{refsection}[bl-author, bl-registered]
        % Это refsection=2.
        % Процитированные здесь работы:
        %  * попадают в авторскую библиографию, при usefootcite==0 и стиле `\insertbiblioauthorimportant`.
        %  * ни на что не влияют в противном случае
        \nocite{sitito_2}
        \nocite{vakbib2}%vak
        \nocite{patbib1}%patent
        \nocite{progbib1}%program
        \nocite{bib1}%other
        \nocite{confbib1}%conf
    \end{refsection}%
        %
        % Всё, что вне этих двух refsection, это refsection=0,
        %  * для диссертации - это нормальные ссылки, попадающие в обычную библиографию
        %  * для автореферата:
        %     * при usefootcite==0, ссылка корректно сработает только для источника из `external.bib`. Для своих работ --- напечатает "[0]" (и даже Warning не вылезет).
        %     * при usefootcite==1, ссылка сработает нормально. В авторской библиографии будут только процитированные в refsection=0 работы.
}
